\section{Experiments and Results}

nnU-Net's automatic configuration search selected the softmax-averaged ensemble of
the 2D and 3D full-resolution configurations for all three label regimes, yielding
ten trained model-folds per regime (thirty in total). The four subsections below
report results for each experiment defined in ``\nameref{sec:design}''
(\autoref{tab:experiments}).

\subsection{Global Fairness Audit}
\label{sec:global}

We first evaluate $M_{\mathrm{mix}}$ on the full $228$-case test set
$\mathcal{D}_{\mathrm{test}}$, scoring each exam against whichever reference label
it carries.

\paragraph{Segmentation quality.}
$M_{\mathrm{mix}}$ achieves mean Dice of $0.959$ (vertebral body), $0.935$ (disc),
and $0.947$ (macro average) on the $228$ test cases. The median HD95 is
$0.34$\,mm (macro), and nDSC tracks Dice closely (macro $0.949$), confirming that
the volume correction of ``\nameref{sec:confounders}'' does not materially change
scores in this cohort. Restricting to the $76$ gold-labelled test cases -- the only
ones scored against expert references -- gives a macro Dice of $0.897$, comparable
to the $0.916$ Zhou et al.~\cite{zhou2025cspineseg} report on their $100$-case gold
test set with a different split and smaller training pool ($391$ expert-annotated
cases versus our $798$ mixed).

\paragraph{Fairness verdict.}
\autoref{tab:global_dir} reports the disparate impact ratio for the three attributes
in $\mathcal{A}$ across all three metrics.
Every DIR exceeds $0.96$ and all lower confidence bounds sit above the four-fifths
threshold of $0.80$ -- no grouping approaches adverse impact on any metric. Across
all seven demographic groupings and nine score columns ($63$ tests), zero reach
significance after Benjamini--Hochberg correction (smallest
$p_{\mathrm{fdr}} = 0.72$). Effect sizes are negligible throughout: the
rank-biserial correlation $|r_{rb}| \le 0.14$ for sex and race on the macro
scores, and $\varepsilon^2 < 0.01$ for age. A multivariable OLS regression of macro Dice on
sex, race, and age jointly yields $R^2 = 0.013$ ($F$-test $p = 0.43$), with no
covariate reaching significance (sex $p = 0.22$, race $p = 0.94$, age $p = 0.24$)
-- the protected attributes collectively explain $1.3\%$ of performance variance.
The sensitivity sweep confirms that the verdict is robust across a range of
beneficial-outcome thresholds. In short, $M_{\mathrm{mix}}$ is fair across
all three protected attributes: no demographic group is disadvantaged by any
practical measure.

\autoref{fig:global_audit} (left) visualises the DIRs with their bootstrap confidence
intervals: all estimates cluster near $1.0$, with every lower bound clearing
$0.80$. \autoref{fig:global_audit} (right) shows the score distribution that later
experiments are most sensitive to -- macro Dice by age bin. The three groups have
nearly identical medians (${\sim}0.974$), but the youngest bin ($<40$) carries a
fat lower tail driven by a handful of difficult cases; the age gradient itself
is negligible.

\begin{table}[t]
\centering
\caption{Disparate impact ratio (DIR) and demographic parity difference (DPD) with
BCa bootstrap $95\%$ confidence intervals on DIR, for $M_{\mathrm{mix}}$ on the
full $228$-case test set. DIR $= 1.0$ and DPD $= 0$ denote parity; the four-fifths
threshold ($0.80$) is the conventional adverse-impact line. All DIR lower bounds
exceed $0.80$.}
\label{tab:global_dir}
\small
\begin{tabular*}{\textwidth}{@{\extracolsep{\fill}} l c c c c c c}
\hline\noalign{\smallskip}
 & \multicolumn{2}{c}{Dice} & \multicolumn{2}{c}{nDSC} & \multicolumn{2}{c}{HD95} \\
\noalign{\smallskip}\cline{2-3}\cline{4-5}\cline{6-7}\noalign{\smallskip}
Grouping & DIR & DPD & DIR & DPD & DIR & DPD \\
\noalign{\smallskip}\hline\noalign{\smallskip}
Sex          & $0.991\;[0.97,\,1.00]$ & $0.009$ & $0.983\;[0.94,\,1.00]$ & $0.018$ & $0.991\;[0.97,\,1.00]$ & $0.009$ \\
Race (W/B)   & $0.976\;[0.90,\,1.00]$ & $0.024$ & $0.969\;[0.89,\,1.00]$ & $0.031$ & $0.970\;[0.89,\,1.00]$ & $0.028$ \\
Age (3 bins) & $0.988\;[0.98,\,1.00]$ & $0.012$ & $0.978\;[0.92,\,1.00]$ & $0.022$ & $0.979\;[0.96,\,1.00]$ & $0.020$ \\
\noalign{\smallskip}\hline
\end{tabular*}
\end{table}

\begin{figure}[t]
\centering
\includegraphics[width=0.525\textwidth]{dir_forest_global.png}\hfill
\includegraphics[width=0.455\textwidth]{violin_dice_macro_age.png}
\caption{Global fairness audit of $M_{\mathrm{mix}}$ on the full $228$-case test
set. \textbf{Left:} Forest plot of disparate impact ratios across all seven
demographic groupings and three metrics; point estimates (dots) with BCa $95\%$
bootstrap confidence intervals (whiskers). The dashed line marks the four-fifths
adverse-impact threshold ($0.80$); every lower bound clears it comfortably.
\textbf{Right:} Distribution of macro Dice by age bin ($n = 45$, $86$, $97$).
Medians are nearly identical (${\sim}0.974$); the youngest bin carries a wider
lower tail (a handful of difficult cases), but the overall age effect is negligible
($\varepsilon^2 < 0.01$, $p_{\mathrm{fdr}} = 0.98$).}
\label{fig:global_audit}
\end{figure}

\subsection{Biased Ruler Effect}
\label{sec:biased_ruler}

We now ask whether the choice of reference label can change a fairness verdict.
The global audit of ``\nameref{sec:global}'' found $M_{\mathrm{mix}}$ fair across
all three attributes ($0/63$ FDR-significant tests, all DIRs $> 0.96$) -- but that
verdict rests on expert labels. What happens if we replace the ruler?
We take the same model, the same $76$ gold-test images, and
score them twice: once against the expert (gold) labels -- the ground truth --
and once against $M_{\mathrm{gold}}$'s predictions on those same images, which
serve as the generated-silver ruler described in ``\nameref{sec:design}''.
Because the predictions are identical and only the reference label changes, any
difference in the fairness verdict is a pure ruler effect.

\paragraph{Performance gap between rulers.}
Against the gold ruler, $M_{\mathrm{mix}}$ achieves a macro Dice of $0.897$
(vertebral body $0.922$, disc $0.872$). Against the silver ruler the same
predictions score $0.973$ (vertebral body $0.979$, disc $0.968$) -- an inflation
of nearly $8$ percentage points. The gap itself is a first signal: the silver
ruler is far more lenient than the expert, because it shares systematic errors
with the model under test.

\paragraph{The verdict flip.}
Under the gold ruler, zero of the $63$ statistical tests reach significance after
Benjamini--Hochberg correction (smallest $p_{\mathrm{fdr}} = 0.13$) -- the model
is judged fair. Under the silver ruler, $11$ of $63$ tests are FDR-significant at
$\alpha = 0.05$, and every one of them is an age comparison (the three-bin and
median-split groupings, across Dice and nDSC for both structures and the macro
score). HD95 tells the same inflation story (gold mean $2.50$\,mm vs silver
$0.40$\,mm) but contributes no significant tests on either ruler, because its
heavy-tailed variance overwhelms the age signal on both.
The significance verdict reverses purely from the choice of reference
label. \autoref{tab:biased_ruler} summarises the two rulers side by side.

\begin{table}[t]
\centering
\caption{Biased-ruler comparison on the $76$ gold-test images (macro Dice).
The same predictions of $M_{\mathrm{mix}}$ are scored against two rulers: expert
labels (gold) and $M_{\mathrm{gold}}$'s predictions (silver). DIR and DPD are
rate-based (threshold $\tau = 0.8$); $p_{\mathrm{fdr}}$ is the Kruskal--Wallis
test after Benjamini--Hochberg correction. The silver ruler saturates
(DIR~$\equiv 1.0$, zero failures above $\tau$) while simultaneously producing
significance on the continuous test for age.}
\label{tab:biased_ruler}
\small
\begin{tabular*}{\textwidth}{@{\extracolsep{\fill}} l c c c c c c}
\hline\noalign{\smallskip}
 & \multicolumn{3}{c}{Gold ruler} & \multicolumn{3}{c}{Silver ruler} \\
\noalign{\smallskip}\cline{2-4}\cline{5-7}\noalign{\smallskip}
Grouping & DIR & DPD & $p_{\mathrm{fdr}}$ & DIR & DPD & $p_{\mathrm{fdr}}$ \\
\noalign{\smallskip}\hline\noalign{\smallskip}
Sex          & $1.000$ & $0.000$ & $0.997$ & $1.000$ & $0.000$ & $0.718$ \\
Race (W/B)   & $0.956$ & $0.044$ & $0.997$ & $1.000$ & $0.000$ & $0.328$ \\
Age (3 bins) & $0.952$ & $0.048$ & $0.270$ & $1.000$ & $0.000$ & $0.026^{\ast}$ \\
\noalign{\smallskip}\hline
\end{tabular*}

\smallskip
{\footnotesize $^{\ast}$\,Significant after FDR correction at $\alpha = 0.05$.
All $11$ significant tests on the silver ruler are age groupings (three-bin and
median split, Dice and nDSC).}
\end{table}

\paragraph{Mechanism: false confidence via variance collapse.}
The verdict flip does not arise because the silver ruler reveals a larger gap --
it does the opposite. The age trend is smaller against silver than against
gold: median macro Dice drops from $0.925$ ($<$40) to $0.898$ (60+) on the gold
ruler (a $2.7$-point spread) but only from $0.980$ to $0.974$ on silver (a
$0.6$-point spread). What changes is the within-group variance. Recall that the
silver ruler is $M_{\mathrm{gold}}$'s predictions -- so scoring against it
measures inter-model agreement rather than accuracy. That agreement sits at
${\sim}0.97$ Dice, far above the ${\sim}0.89$ that the same predictions achieve
against the expert labels, and the standard deviation across the $76$ cases drops
from $0.058$ (gold) to $0.014$ (silver) -- a ${\sim}4\times$ reduction. We attribute
this high agreement to the inherent simplicity of cervical-spine segmentation --
vertebral bodies and discs are large, high-contrast structures with well-defined
boundaries -- combined with the shared architecture and training recipe, which
lead both models to converge on similar solutions regardless of the specific
training cases.\footnote{$M_{\mathrm{mix}}$ and $M_{\mathrm{gold}}$ do share $288$
training cases, but this is not the primary driver: even
$M_{\mathrm{gold}}$ and $M_{\mathrm{silver}}$, trained on fully disjoint data
(zero case overlap), still agree at $0.97$ macro Dice on the same $76$ images.
The high inter-model agreement is a property of the task, not of shared
training data.} The Kruskal--Wallis test's power
scales with signal-to-noise, so the same clinically negligible age gradient that
is buried in gold's noise ($H = 7.0$, $p = 0.030$,
$p_{\mathrm{fdr}} = 0.27$) clears correction once silver's collapsed variance
sharpens the ranks ($H = 11.4$, $p = 0.003$, $p_{\mathrm{fdr}} = 0.026$). The
ruler did not change the patient; it changed the test's power.
A threshold sensitivity sweep confirms the saturation: the silver DIR remains
pinned at $1.0$ for all groupings through $\tau = 0.85$ and only drops (to
$0.97$) at $\tau = 0.90$, whereas the gold DIR already flags age disparity at
$\tau = 0.70$ and collapses to $0.57$ at $\tau = 0.90$.
\autoref{fig:biased_ruler} visualises the variance collapse: the same downward
staircase appears on both rulers, but gold's wide boxes bury the signal while
silver's razor-thin boxes expose it.

\paragraph{Comparison with Parikh et al.}
Parikh et al.~\cite{parikh2025labelbias} show a biased-ruler effect in
breast-MRI segmentation where the silver ruler makes a real gap look larger
than it truly is (DIR drops from $0.871$ to $0.815$). Our silver ruler does
something different: it leaves the gap the same size (or even shrinks it) but
makes a negligible gap look statistically significant. In their case, the ruler
exaggerates the signal; in ours, it removes the noise that would otherwise mask
it. Both lead to the same practical problem -- the fairness verdict depends on
which reference label you choose -- but through opposite routes.

\paragraph{The age effect is intrinsic.}
Although the significance is a ruler artifact, the underlying age trend is not.
It appears on both rulers in the same direction (60+ worst), is
consistent with the degenerative anatomy documented in
``\nameref{sec:confounders}'' (vertebral merging in older spines), and is
independently corroborated by the encoder probe
(``\nameref{sec:probing_results}''). An effect visible through the true ruler, the
correlated ruler, and the learned representation is unlikely to be a label
artifact -- it reflects intrinsic difficulty in older spines.

\paragraph{Not unfairness.}
A statistically significant gap is not the same as an unfair model. The
magnitudes here are clinically negligible: the worst-vs-best age group spans
$2.7$ Dice points on gold ($0.925 \to 0.898$) and $0.6$ points on silver, with
every patient receiving a ${\sim}0.90$+ Dice segmentation. All gold-ruler DIRs
exceed $0.94$ -- far above the four-fifths threshold -- and no subgroup is failed
by any practical measure. The contribution of this experiment is methodological:
it demonstrates that a machine-generated reference label can manufacture a
finding out of a non-issue, and that significance alone should not be read as
evidence of harm.

\begin{figure}[t]
\centering
\includegraphics[width=\textwidth]{age_trend_gold_vs_silver.png}
\caption{The variance-collapse mechanism behind the verdict flip. Both panels
show macro Dice by age bin for the same $76$ predictions of $M_{\mathrm{mix}}$,
scored against the gold ruler (left) and the silver ruler (right). The age
gradient runs in the same direction on both (60+ worst), but the gold ruler's
per-group standard deviations ($0.03$--$0.10$) bury the signal below FDR
correction, while the silver ruler's collapsed variance ($0.01$--$0.02$) exposes
it. The ruler changed the test's power, not the patient's outcome.}
\label{fig:biased_ruler}
\end{figure}

\subsection{Bias Amplification}
\label{sec:amplification}

The biased-ruler experiment showed that the choice of reference label can distort
a fairness verdict at evaluation time. We now ask whether label provenance in
training can do the same. E3 and E4 (\autoref{tab:experiments}) compare
the three models -- $M_{\mathrm{mix}}$, $M_{\mathrm{gold}}$, and
$M_{\mathrm{silver}}$ -- on the same $76$ gold-test cases scored against expert
labels, so any difference in fairness is attributable to the training labels
alone.

\paragraph{Segmentation quality.}
$M_{\mathrm{silver}}$ achieves a macro Dice of $0.897$ (vertebral body $0.921$,
disc $0.872$) -- identical to $M_{\mathrm{mix}}$ to three decimal places despite
never seeing an expert label during training. $M_{\mathrm{gold}}$ scores slightly
lower ($0.888$; vertebral body $0.914$, disc $0.862$), likely because it trains on
fewer cases ($288$ versus $450$ and $798$). Silver training labels do not degrade
absolute performance in this setting; the larger training pool compensates for the
noisier supervision.

\paragraph{Fairness verdict.}
\autoref{tab:amplification} reports the macro Dice DIR for all three models. Zero
of the $63$ statistical tests reach significance after Benjamini--Hochberg
correction on any model (smallest $p_{\mathrm{fdr}} = 0.061$,
$M_{\mathrm{gold}}$ nDSC disc by age). The
$M_{\mathrm{silver}}$ DIRs are nearly identical to $M_{\mathrm{mix}}$ -- in most
groupings the values coincide exactly -- and all exceed $0.93$. If anything,
$M_{\mathrm{gold}}$ is the least fair of the three on the disc (race W/B/O DIR
$= 0.813$ versus $0.875$ for both $M_{\mathrm{mix}}$ and $M_{\mathrm{silver}}$),
though even that value clears the four-fifths threshold. There is no evidence that
silver training labels amplify demographic bias.

\begin{table}[t]
\centering
\caption{Bias amplification: macro Dice DIR for the three models on the $76$
gold-test cases (expert reference). DIR $= 1.0$ denotes parity; the four-fifths
threshold is $0.80$. $M_{\mathrm{silver}}$ matches $M_{\mathrm{mix}}$ almost
exactly; $M_{\mathrm{gold}}$ is occasionally worse. All values exceed
$0.80$; $0/63$ tests reach FDR significance on any model.}
\label{tab:amplification}
\small
\begin{tabular*}{\textwidth}{@{\extracolsep{\fill}} l c c c}
\hline\noalign{\smallskip}
Grouping & $M_{\mathrm{mix}}$ & $M_{\mathrm{gold}}$ & $M_{\mathrm{silver}}$ \\
\noalign{\smallskip}\hline\noalign{\smallskip}
Sex          & $1.000$ & $0.946$ & $1.000$ \\
Race (W/B)   & $0.956$ & $0.994$ & $0.956$ \\
Race (W/B/O) & $0.938$ & $0.938$ & $0.938$ \\
Age (3 bins) & $0.952$ & $0.956$ & $0.952$ \\
Age (median) & $0.998$ & $0.996$ & $0.998$ \\
\noalign{\smallskip}\hline
\end{tabular*}
\end{table}

\paragraph{Interpretation.}
The null result is expected: bias amplification requires both an underlying
disparity for noisy labels to exaggerate and a silver source that is
differentially biased across groups. Neither precondition holds here. The global
audit (``\nameref{sec:global}'') already found no demographic gap to amplify
($0/63$, all DIRs $> 0.96$), and the silver labels are high-quality -- Zhou et
al.~\cite{zhou2025cspineseg} selected which exams to expert-annotate
pseudo-randomly (by medical record number), leaving the gold and silver pools
demographically comparable, so the machine-generated labels carry no systematic
group-level bias for training to propagate.
This contrasts with Parikh et al.~\cite{parikh2025labelbias}, who observed $66\%$
gap widening (DIR dropping below $0.80$) when training on their silver labels. The
difference lies in the preconditions: their silver annotator was an independently
trained model that was itself biased against younger patients, acting on a task
with a substantial pre-existing age disparity. In our setting, the silver labels
are produced by a model trained on a representative subset of the same
institution's gold annotations, segmenting large high-contrast structures whose
boundaries may leave little room for demographically differential error --
though we cannot test this hypothesis directly. The result suggests that
mixing silver labels into training is unlikely to harm fairness when the
underlying labels are well-calibrated and no pre-existing disparity exists for
noisy supervision to exaggerate.

\subsection{Demographic Probing}
\label{sec:probing_results}

The preceding experiments ask whether demographic groups receive different
outputs. We now ask a logically independent question: does the model's
internal representation carry demographic signal at all? Encoding a protected
attribute and producing a performance disparity are separate phenomena -- a model
can do one without the other~\cite{petersen2023feature}. The probe setup --
two encoders (nnU-Net bottleneck and MRI-CORE), each paired with a random-init
null, scored via PCA(50) and logistic regression in 5-fold CV -- is defined in
``\nameref{sec:probing_method}'' (\autoref{sec:probing_method}).

\paragraph{Results.}
\autoref{tab:probe} reports probe scores for all four encoder--null pairs across
three demographic attributes plus manufacturer (a sanity check). Sex is the most
strongly encoded attribute: the nnU-Net bottleneck reaches $0.957$ AUROC, with the
segmentation task contributing $+0.186$ above its random null ($0.771$). MRI-CORE
scores $0.931$, but its random-init null already achieves $0.836$ -- indicating
that ${\sim}90\%$ of the sex signal in the task-agnostic encoder comes from raw
pixel statistics (body silhouette, FOV extent) rather than learned anatomy. Age
is moderately encoded on both ($0.642$ nnU-Net, $0.632$ MRI-CORE), with task
training again the larger contributor ($+0.166$ vs.\ $+0.098$).
Race sits near chance on all encoders (wide CIs, fold variance $0.41$--$0.55$), consistent
with the null race finding in the fairness audit. Manufacturer reaches ${\sim}0.997$
universally -- a pixel-level scanner fingerprint independent of learned weights,
serving as a healthy end-to-end sanity check.

A confounder analysis on the full cohort confirms that the probe signals are not
artefacts of scanner assignment or attribute entanglement: sex $\times$
manufacturer ($\chi^2$ $p = 0.21$, Cram\'er's $V = 0.035$) and age $\times$ sex
($p = 0.08$, $r_{rb} = 0.07$) are not significant; age $\times$ manufacturer
reaches significance ($p = 0.02$) but at a negligible effect size
($r_{rb} = -0.09$), meaning older patients are only marginally more likely to
appear on one scanner -- too weak an association to account for the age probe
scores observed above.

\begin{table}[t]
\centering
\caption{Linear-probe scores for demographic attributes on frozen encoder
features ($n = 1{,}085$). ``Null'' denotes a random-initialised network of
identical architecture and preprocessing. $\Delta_{\mathrm{task}}$ is the
trained-minus-null gap, isolating the contribution of the training objective.
Sex is strongly encoded yet produces no performance disparity
(``\nameref{sec:global}''); age is the one axis with both encoding and a
(negligible) performance trend. Values are mean $\pm$ Nadeau--Bengio 95\% CI.}
\label{tab:probe}
\small
\begin{tabular*}{\textwidth}{@{\extracolsep{\fill}} l l c c c}
\hline\noalign{\smallskip}
Encoder & Attribute & Trained & Null & $\Delta_{\mathrm{task}}$ \\
\noalign{\smallskip}\hline\noalign{\smallskip}
nnU-Net        & Sex (AUROC)       & $0.957\pm.030$ & $0.771\pm.066$ & $+0.186$ \\
(segmentation) & Age (bal.\ acc.)  & $0.642\pm.023$ & $0.476\pm.055$ & $+0.166$ \\
               & Race (bal.\ acc.) & $0.491\pm.104$          & $0.373\pm.052$ & $+0.118$ \\
               & Manuf.\ (AUROC)   & $0.997\pm.004$          & $0.983\pm.012$ & $+0.014$ \\
\noalign{\smallskip}\cline{1-5}\noalign{\smallskip}
MRI-CORE       & Sex (AUROC)       & $0.931\pm.031$ & $0.836\pm.039$ & $+0.095$ \\
(task-agnostic) & Age (bal.\ acc.) & $0.632\pm.043$ & $0.534\pm.083$ & $+0.098$ \\
               & Race (bal.\ acc.) & $0.513\pm.071$ & $0.426\pm.064$ & $+0.087$ \\
               & Manuf.\ (AUROC)   & $0.997\pm.007$ & $0.992\pm.013$ & $+0.005$ \\
\noalign{\smallskip}\hline
\end{tabular*}
\end{table}

\paragraph{Encoding does not imply disparity.}
The central finding is a dissociation: sex is by far the most strongly encoded
attribute ($0.957$), yet the global audit found zero sex disparity ($0/63$
significant tests, DIR~$= 0.991$). The segmentation task has learned to exploit
sex-linked vertebral morphology -- larger bodies and discs in men, different canal
geometry -- without producing differential outcomes. Conversely, age is only
moderately encoded ($0.642$) yet is the one axis that shows a (clinically
negligible) performance trend, corroborating the biased-ruler finding that the age
effect is intrinsic anatomical difficulty rather than a label artefact. The probe
thus confirms that encoding is neither necessary nor sufficient for disparity: a
model can strongly represent a demographic feature and still be fair, because the
decoder maps those different representations to equally accurate outputs.